\chapter{The main theorem}

\section{Statement}\label{sec:statement}

\begin{theorem}[Main]
\label{thm:main}
\lean{Example.main,
      Example.main'}
\uses{def:poly, lem:degree}
For every $n$ the identities
\begin{align}
  a &= b \\
  b &= c
\end{align}
hold, and so does
\begin{equation}
  \label{eq:gauss}
  \sum_{i=0}^{n} i = \frac{n(n+1)}{2}.
\end{equation}
See \Cref{lem:degree} and \cref{def:poly,lem:missing}.
\end{theorem}

\begin{proof}
\uses{lem:degree, lem:nowhere}
\leanok
By induction.  The steps are
\begin{enumerate}
  \item the base case, and
  \item the inductive step, which uses
    \begin{itemize}
      \item the degree lemma, and
      \item nothing else.
    \end{itemize}
\end{enumerate}
\end{proof}

\begin{openproblem}[Still open]
\label{op:open}
\statusnote{Nobody has tried.}
Is $\Fq$ algebraically closed?  Obviously not.
\end{openproblem}
